\documentclass[UTF8,12pt,a4paper]{ctexart}
\usepackage{geometry}
\usepackage{longtable}
\usepackage{array}
\usepackage{booktabs}
\usepackage{hyperref}
\usepackage{xcolor}
\usepackage{enumitem}
\usepackage{amsmath}
\usepackage{amssymb}
\usepackage{listings}
\geometry{left=2.5cm,right=2.5cm,top=2.6cm,bottom=2.6cm}
\hypersetup{colorlinks=true,linkcolor=blue,urlcolor=blue}
\lstset{basicstyle=\ttfamily\small,breaklines=true,frame=single,columns=fullflexible}

\title{RDDQF v1.2 平台功能增强 Agent 修改指南\\\large 训练数据消费、批次抽检判定、质检任务池管理与状态溯源}
\author{Robot Demonstration Data Quality Framework}
\date{v1.2}

\begin{document}
\maketitle
\tableofcontents
\newpage

\section{目标与范围}

本次修改目标是将现有机器人数据质检平台从“质检结果记录系统”进一步升级为“训练数据消费与数据资产供给系统”。

当前平台已基本完成：

\begin{center}
上传 $\rightarrow$ MinIO 存储 $\rightarrow$ 抽检质检 $\rightarrow$ 记录质检结果
\end{center}

本次需要补齐后续链路：

\begin{center}
质检结果 $\rightarrow$ 批次判定 $\rightarrow$ Episode 最终训练可用状态 $\rightarrow$ 任务级可用数据统计 $\rightarrow$ 合格数据清单导出
\end{center}

本指南要求 Agent 完成以下四类功能：

\begin{enumerate}
  \item 训练数据集管理页面增强；
  \item 批次抽检驳回逻辑与 Episode 最终状态写入；
  \item 管理员对质检员任务池的管理能力；
  \item 数据导出字段增强与状态溯源展示。
\end{enumerate}

\section{核心业务规则确认}

\subsection{批次失败率计算口径}

本版本明确采用“按抽检数计算”的失败率逻辑，而不是按批次总数计算。

设：

\begin{itemize}
  \item $N_{sampled}$：该批次已抽检并完成质检的 Episode 数量；
  \item $N_{manual\_fail}$：该批次中人工质检判定为不合格的 Episode 数量；
  \item $\theta$：批次驳回阈值，默认 $\theta=0.10$。
\end{itemize}

则批次失败率为：

\[
FailureRate=\frac{N_{manual\_fail}}{N_{sampled}}
\]

当：

\[
FailureRate>\theta
\]

则整批数据被驳回。

\textbf{注意：}失败率分母必须为已完成抽检的样本数，而不是批次总 Episode 数。

\subsection{批次状态}

批次状态建议使用：

\begin{lstlisting}
batch_decision:
  - pending
  - accepted
  - rejected
\end{lstlisting}

其中：

\begin{itemize}
  \item \texttt{pending}：批次尚未完成足够的质检判断；
  \item \texttt{accepted}：批次抽检失败率未超过阈值；
  \item \texttt{rejected}：批次抽检失败率超过阈值，整批驳回。
\end{itemize}

\section{Episode 最终状态模型}

现有系统不能只使用单一的 \texttt{pass/fail} 字段。必须区分人工判定与批次规则推导判定。

\subsection{推荐字段}

在 Episode 模型或 Episode 关联状态表中增加以下字段：

\begin{longtable}{p{4.5cm}p{4cm}p{6cm}}
\toprule
字段 & 类型/枚举 & 含义 \\
\midrule
\texttt{manual\_qc\_status} & enum & 人工质检结果 \\
\texttt{final\_dataset\_status} & enum & 下游训练是否可用 \\
\texttt{final\_decision\_source} & enum & 最终状态来源 \\
\texttt{final\_decision\_reason} & text & 最终判定原因 \\
\texttt{final\_decided\_at} & datetime & 最终状态更新时间 \\
\texttt{is\_exportable} & bool & 是否可导出给训练组 \\
\bottomrule
\end{longtable}

\subsection{人工质检状态}

\begin{lstlisting}
manual_qc_status:
  - not_reviewed
  - manual_pass
  - manual_fail
\end{lstlisting}

\subsection{最终训练可用状态}

\begin{lstlisting}
final_dataset_status:
  - pending
  - qualified
  - unqualified
\end{lstlisting}

\subsection{最终状态来源}

\begin{lstlisting}
final_decision_source:
  - manual_pass
  - manual_fail
  - batch_accept_inferred_pass
  - batch_reject_propagated_fail
  - batch_reject_override_manual_pass
  - pending
\end{lstlisting}

含义如下：

\begin{longtable}{p{5cm}p{10cm}}
\toprule
来源 & 含义 \\
\midrule
\texttt{manual\_pass} & 该 Episode 被质检员人工判定为合格，且所在批次未被驳回 \\
\texttt{manual\_fail} & 该 Episode 被质检员人工判定为不合格 \\
\texttt{batch\_accept\_inferred\_pass} & 所在批次通过抽检，未被抽检的 Episode 被推断为合格 \\
\texttt{batch\_reject\_propagated\_fail} & 所在批次被驳回，未抽检 Episode 被连带判定为不合格 \\
\texttt{batch\_reject\_override\_manual\_pass} & 该 Episode 人工判定为合格，但所在批次整体被驳回，因此最终不可用 \\
\texttt{pending} & 尚未产生最终状态 \\
\bottomrule
\end{longtable}

\section{批次判定服务}

\subsection{新增服务}

建议新增或完善：

\begin{lstlisting}
app/services/batch_adjudication.py
\end{lstlisting}

核心类：

\begin{lstlisting}
class BatchAdjudicationService:
    def adjudicate_batch(batch_id: int) -> BatchDecisionResult:
        ...

    def recompute_batch_decision(batch_id: int, reason: str) -> BatchDecisionResult:
        ...
\end{lstlisting}

\subsection{触发时机}

每次质检员提交 episode QC 结果后，系统应检查该 Episode 所属 batch 是否满足批次判定条件。

推荐触发链路：

\begin{center}
提交人工质检结果 $\rightarrow$ 更新 \texttt{manual\_qc\_status} $\rightarrow$ 调用 \texttt{adjudicate\_batch(batch\_id)} $\rightarrow$ 写入 batch 决策 $\rightarrow$ 更新 episode 最终状态
\end{center}

\subsection{批次判定条件}

若该批次已分配抽检任务，则只有当所有抽检任务完成后，才进行批次判定。

设：

\begin{itemize}
  \item $N_{sampled}$：被抽检的 Episode 数；
  \item $N_{reviewed}$：已完成质检的抽检 Episode 数。
\end{itemize}

只有当：

\[
N_{reviewed}=N_{sampled}
\]

时执行批次判定。否则保持：

\begin{lstlisting}
batch_decision = pending
\end{lstlisting}

\subsection{批次判定算法}

\begin{lstlisting}
def adjudicate_batch(batch_id):
    total_count = count_all_episodes(batch_id)
    sampled_count = count_sampled_episodes(batch_id)
    reviewed_count = count_reviewed_sampled_episodes(batch_id)
    manual_fail_count = count_manual_failed_sampled_episodes(batch_id)

    if sampled_count == 0 or reviewed_count < sampled_count:
        batch_decision = 'pending'
        return

    failure_rate = manual_fail_count / sampled_count

    if failure_rate > reject_threshold:
        batch_decision = 'rejected'
    else:
        batch_decision = 'accepted'

    update_batch_decision(...)
    update_episode_final_statuses(batch_id, batch_decision)
    create_batch_decision_log(...)
\end{lstlisting}

\section{Episode 最终状态更新规则}

\subsection{批次未完成判定}

若：

\begin{lstlisting}
batch_decision = pending
\end{lstlisting}

则：

\begin{lstlisting}
final_dataset_status = pending
final_decision_source = pending
is_exportable = False
\end{lstlisting}

\subsection{批次通过}

若：

\begin{lstlisting}
batch_decision = accepted
\end{lstlisting}

则：

\begin{longtable}{p{4cm}p{4cm}p{5cm}p{2cm}}
\toprule
人工状态 & 最终状态 & 来源 & 可导出 \\
\midrule
\texttt{manual\_pass} & \texttt{qualified} & \texttt{manual\_pass} & true \\
\texttt{manual\_fail} & \texttt{unqualified} & \texttt{manual\_fail} & false \\
\texttt{not\_reviewed} & \texttt{qualified} & \texttt{batch\_accept\_inferred\_pass} & true \\
\bottomrule
\end{longtable}

\subsection{批次驳回}

若：

\begin{lstlisting}
batch_decision = rejected
\end{lstlisting}

则：

\begin{longtable}{p{4cm}p{4cm}p{5cm}p{2cm}}
\toprule
人工状态 & 最终状态 & 来源 & 可导出 \\
\midrule
\texttt{manual\_pass} & \texttt{unqualified} & \texttt{batch\_reject\_override\_manual\_pass} & false \\
\texttt{manual\_fail} & \texttt{unqualified} & \texttt{manual\_fail} & false \\
\texttt{not\_reviewed} & \texttt{unqualified} & \texttt{batch\_reject\_propagated\_fail} & false \\
\bottomrule
\end{longtable}

\section{批次判定日志}

新增表或模型：

\begin{lstlisting}
BatchDecisionLog
\end{lstlisting}

推荐字段：

\begin{longtable}{p{4.5cm}p{10cm}}
\toprule
字段 & 含义 \\
\midrule
\texttt{id} & 日志 ID \\
\texttt{batch\_id} & 批次 ID \\
\texttt{policy\_version} & 判定策略版本，例如 \texttt{batch-reject-v1.2} \\
\texttt{reject\_threshold} & 驳回阈值，默认 0.10 \\
\texttt{failure\_rate\_denominator} & 固定为 \texttt{sampled\_count} \\
\texttt{total\_episode\_count} & 批次总 Episode 数 \\
\texttt{sampled\_episode\_count} & 抽检 Episode 数 \\
\texttt{reviewed\_episode\_count} & 已完成质检数 \\
\texttt{manual\_fail\_count} & 人工不合格数 \\
\texttt{failure\_rate} & 失败率 \\
\texttt{batch\_decision} & 批次判定结果 \\
\texttt{created\_at} & 日志创建时间 \\
\texttt{created\_by} & 操作者，自动判定则为 system \\
\texttt{reason} & 判定原因 \\
\bottomrule
\end{longtable}

\section{训练数据集管理页面增强}

\subsection{页面定位}

新增或增强页面：

\begin{lstlisting}
/dataset-management
\end{lstlisting}

页面面向：

\begin{itemize}
  \item 下游模型训练人员；
  \item 数据采集团队；
  \item 质检管理员。
\end{itemize}

核心问题：

\begin{quote}
某个任务当前有多少条可用于训练的合格 Episode？是否达到下游需求数量？能否导出合格数据清单？
\end{quote}

\subsection{任务级统计卡片}

选择任务类型后，顶部显示：

\begin{lstlisting}
任务名称：切土豆
目标数量：300
当前合格：193
缺口数量：107
完成率：64.3%
已采集批次：3
已通过批次：2
已驳回批次：1
\end{lstlisting}

合格数量计算：

\[
QualifiedCount=Count(final\_dataset\_status=qualified)
\]

缺口数量：

\[
Gap=max(TargetCount-QualifiedCount,0)
\]

\subsection{批次列表}

批次表格字段建议：

\begin{longtable}{p{2.5cm}p{2cm}p{2cm}p{2cm}p{2cm}p{2.2cm}p{2cm}}
\toprule
批次 & 总数 & 抽检数 & 人工失败 & 失败率 & 批次状态 & 可用数 \\
\midrule
batch-001 & 100 & 25 & 4 & 16\% & rejected & 0 \\
batch-002 & 100 & 25 & 2 & 8\% & accepted & 98 \\
\bottomrule
\end{longtable}

失败率必须显示为：

\[
\frac{ManualFailCount}{SampledCount}
\]

前端文案应为：

\begin{quote}
失败率 = 人工不合格数 / 抽检完成数
\end{quote}

\subsection{Episode 列表}

默认展示：

\begin{lstlisting}
final_dataset_status = qualified
\end{lstlisting}

支持筛选：

\begin{itemize}
  \item 任务类型；
  \item 批次；
  \item 最终状态；
  \item 最终判定来源；
  \item 人工质检状态；
  \item L3 分数区间；
  \item 创建时间；
  \item 质检完成时间。
\end{itemize}

\section{导出字段增强}

\subsection{导出原则}

导出对象为 Episode 级元数据清单，不直接导出 MinIO 文件本体。

导出过滤条件：

\begin{lstlisting}
final_dataset_status = qualified
is_exportable = True
\end{lstlisting}

\subsection{导出格式}

本版本建议支持：

\begin{itemize}
  \item CSV：便于人工查看；
  \item JSON：便于接口消费；
  \item JSONL：便于训练管线逐行读取。
\end{itemize}

\subsection{推荐导出字段}

必须包含：

\begin{longtable}{p{5cm}p{10cm}}
\toprule
字段 & 含义 \\
\midrule
\texttt{episode\_id} & Episode ID \\
\texttt{episode\_uid} & 全局唯一 Episode 标识，如存在 \\
\texttt{task\_type\_id} & 任务类型 ID \\
\texttt{task\_type\_name} & 任务类型名称 \\
\texttt{batch\_id} & 批次 ID \\
\texttt{batch\_name} & 批次名称 \\
\texttt{final\_dataset\_status} & 最终训练可用状态 \\
\texttt{final\_decision\_source} & 最终状态来源 \\
\texttt{final\_decision\_reason} & 判定原因 \\
\texttt{manual\_qc\_status} & 人工质检状态 \\
\texttt{reviewer\_id} & 质检员 ID，如有 \\
\texttt{qc\_result\_id} & 人工质检记录 ID \\
\texttt{l3\_training\_quality\_score} & L3 v2 训练质量分数 \\
\texttt{motion\_quality\_score} & Motion Quality 分数 \\
\texttt{learnability\_score} & Learnability 分数 \\
\texttt{data\_integrity\_score} & Data Integrity 分数 \\
\texttt{execution\_diagnostic\_score} & Execution Diagnostic 分数 \\
\texttt{minio\_raw\_prefix} & raw 数据 MinIO 路径前缀 \\
\texttt{minio\_processed\_prefix} & processed 数据 MinIO 路径前缀 \\
\texttt{telemetry\_path} & telemetry.npz 路径 \\
\texttt{metadata\_path} & metadata.json 路径 \\
\texttt{manifest\_path} & manifest.json 路径 \\
\texttt{video\_paths} & 视频路径 JSON 字符串 \\
\texttt{created\_at} & Episode 创建时间 \\
\texttt{uploaded\_at} & 上传时间，如有 \\
\texttt{qc\_completed\_at} & 质检完成时间 \\
\texttt{final\_decided\_at} & 最终状态判定时间 \\
\bottomrule
\end{longtable}

\subsection{导出历史}

建议增加导出历史表：

\begin{lstlisting}
DatasetExportJob
\end{lstlisting}

字段建议：

\begin{itemize}
  \item \texttt{id}
  \item \texttt{task\_type\_id}
  \item \texttt{export\_format}
  \item \texttt{episode\_count}
  \item \texttt{file\_path}
  \item \texttt{created\_by}
  \item \texttt{created\_at}
  \item \texttt{filters\_json}
\end{itemize}

\section{管理员质检任务池管理}

\subsection{需求背景}

当前任务派发后会进入质检员任务池。如果管理员无法管理已分配任务，则会出现以下问题：

\begin{itemize}
  \item 质检员请假或离线，任务无法回收；
  \item 某质检员任务过多，无法平衡负载；
  \item 错误派发后无法撤销；
  \item 管理员无法处理卡住的任务。
\end{itemize}

因此需要新增 Reviewer Task Manager。

\subsection{入口设计}

在管理员工作台的审核员工作量列表中，每个审核员后增加按钮：

\begin{lstlisting}
管理任务
\end{lstlisting}

点击后打开 Drawer 或 Modal，展示该审核员任务池。

\subsection{任务池表格字段}

\begin{longtable}{p{3cm}p{3cm}p{3cm}p{2.5cm}p{2.5cm}}
\toprule
字段 & 含义 & 示例 & 是否筛选 & 备注 \\
\midrule
任务 ID & QC task id & 123 & 是 & \\
Episode ID & episode id & ep-001 & 是 & \\
任务类型 & task type & 切土豆 & 是 & \\
批次 & batch & batch-001 & 是 & \\
状态 & pending/in\_progress/completed & pending & 是 & \\
派发时间 & assigned at & 2026-xx-xx & 否 & \\
开始时间 & started at & 2026-xx-xx & 否 & \\
耗时 & duration & 8.5min & 否 & \\
\bottomrule
\end{longtable}

\subsection{支持操作}

\begin{enumerate}
  \item 撤回未开始任务；
  \item 转派未开始任务；
  \item 释放未开始任务回公共任务池；
  \item 批量撤回；
  \item 批量转派；
  \item 查看已完成任务，不允许修改。
\end{enumerate}

\subsection{状态限制}

\begin{longtable}{p{4cm}p{8cm}}
\toprule
任务状态 & 允许操作 \\
\midrule
\texttt{pending} & 撤回、转派、释放回公共池 \\
\texttt{in\_progress} & 默认不可操作；如实现强制释放，必须记录原因 \\
\texttt{completed} & 不可撤销，不可转派，仅可查看 \\
\bottomrule
\end{longtable}

\section{任务操作日志}

新增任务操作日志：

\begin{lstlisting}
TaskOperationLog
\end{lstlisting}

字段建议：

\begin{longtable}{p{5cm}p{10cm}}
\toprule
字段 & 含义 \\
\midrule
\texttt{id} & 日志 ID \\
\texttt{task\_id} & 被操作的质检任务 ID \\
\texttt{episode\_id} & Episode ID \\
\texttt{operation} & 操作类型：revoke / reassign / release / force\_release \\
\texttt{from\_reviewer\_id} & 原质检员 \\
\texttt{to\_reviewer\_id} & 新质检员，如有 \\
\texttt{operator\_id} & 管理员 ID \\
\texttt{reason} & 操作原因 \\
\texttt{created\_at} & 操作时间 \\
\bottomrule
\end{longtable}

\section{API 修改建议}

\subsection{Dataset API}

建议新增或增强：

\begin{lstlisting}
GET  /api/dataset/tasks
GET  /api/dataset/tasks/{task_type_id}/summary
GET  /api/dataset/tasks/{task_type_id}/batches
GET  /api/dataset/tasks/{task_type_id}/episodes
POST /api/dataset/tasks/{task_type_id}/exports
GET  /api/dataset/exports/{export_id}/download
GET  /api/dataset/exports
\end{lstlisting}

\subsection{Batch API}

\begin{lstlisting}
POST /api/batches/{batch_id}/adjudicate
POST /api/batches/{batch_id}/recompute-decision
GET  /api/batches/{batch_id}/decision-log
\end{lstlisting}

\subsection{Reviewer Task Manager API}

\begin{lstlisting}
GET  /api/admin/reviewers/{reviewer_id}/tasks
POST /api/admin/qc-tasks/{task_id}/revoke
POST /api/admin/qc-tasks/{task_id}/reassign
POST /api/admin/qc-tasks/{task_id}/release
POST /api/admin/qc-tasks/bulk-reassign
POST /api/admin/qc-tasks/bulk-revoke
GET  /api/admin/qc-task-operations
\end{lstlisting}

\section{前端修改建议}

\subsection{训练数据集管理页面}

页面需要包含：

\begin{enumerate}
  \item 任务选择器；
  \item 目标数量、合格数量、缺口数量、完成率；
  \item 批次汇总表；
  \item Episode 合格数据表；
  \item 导出按钮；
  \item 导出历史记录；
  \item Episode 最终状态溯源面板。
\end{enumerate}

\subsection{管理员任务管理弹窗}

在审核员工作量列表中增加“管理任务”按钮。

弹窗包含：

\begin{itemize}
  \item 当前审核员信息；
  \item 任务池列表；
  \item 状态筛选；
  \item 批量选择；
  \item 撤回、转派、释放操作；
  \item 操作原因输入框；
  \item 操作日志入口。
\end{itemize}

\subsection{Episode 状态溯源面板}

在 Episode 列表点击某条记录时展示：

\begin{lstlisting}
最终状态：qualified / unqualified
最终来源：manual_pass / batch_accept_inferred_pass / ...
人工质检状态：manual_pass / manual_fail / not_reviewed
所在批次状态：accepted / rejected / pending
批次失败率：manual_fail_count / sampled_count
策略版本：batch-reject-v1.2
判定时间：...
判定原因：...
\end{lstlisting}

\section{MQ-02 与 DX-01 后续优化说明}

\subsection{MQ-02 Hand Action Continuity}

当前 MQ-02 如果仅对 Arm Action 做二阶差分，则无法检测 Hand Action 的突变。

后续建议改为：

\[
C_t^{arm}=\left\|\Delta^2 a_t^{arm}\right\|_{RMS}
\]

\[
C_t^{hand}=\left\|\Delta^2 a_t^{hand}\right\|_{RMS}
\]

其中 Hand Action 应归一化到 $[0,1]$。

最终：

\[
C_t=max(C_t^{arm},C_t^{hand})
\]

或：

\[
C_t=0.7C_t^{arm}+0.3C_t^{hand}
\]

本版本可以暂不实现，但应记录为技术债。

\subsection{DX-01 真正 Lag Alignment}

真正的 lag alignment 应计算：

\[
k^*=\arg\min_k P_{50}(\|a_{t-k}-q_t\|)
\]

而不是对已计算好的 error 序列做平移。

推荐后续实现：

\begin{lstlisting}
for k in range(0, max_lag + 1):
    aligned_action = actions[:-k] if k > 0 else actions
    aligned_qpos = qpos[k:] if k > 0 else qpos
    err = rms(aligned_action - aligned_qpos)
    score_k = percentile(err, 50)
select k with minimal score_k
\end{lstlisting}

本版本可暂不实现，但应在代码注释或 issue 中标记。

\section{测试要求}

Agent 完成修改后，应至少验证：

\begin{enumerate}
  \item 抽检 25 条，失败 4 条时，失败率为 $4/25=16\%$；
  \item 阈值 10\% 下，该批次被 rejected；
  \item rejected 批次下所有 Episode 最终不可导出；
  \item accepted 批次下未抽检 Episode 被推断为 qualified；
  \item 导出文件只包含 \texttt{final\_dataset\_status=qualified} 的 Episode；
  \item 导出字段包含 MinIO 路径、L3 分数、状态来源；
  \item 管理员可撤回 pending 任务；
  \item 管理员可转派 pending 任务；
  \item completed 任务不可撤回；
  \item 所有任务池操作均写入 TaskOperationLog。
\end{enumerate}

\section{Agent 修改 Checklist}

\subsection{后端}

\begin{itemize}
  \item 确认批次失败率为 manual fail / sampled count；
  \item 完善 BatchAdjudicationService；
  \item 增加 Episode final status 字段；
  \item 增加 BatchDecisionLog；
  \item 增加 DatasetExportJob；
  \item 增加 TaskOperationLog；
  \item 增强 Dataset Summary API；
  \item 增强 Dataset Export API；
  \item 增加 Reviewer Task Manager API。
\end{itemize}

\subsection{前端}

\begin{itemize}
  \item 增强数据集管理页面；
  \item 增加导出历史；
  \item 增加 Episode 决策溯源面板；
  \item 在审核员工作量卡片后增加“管理任务”；
  \item 实现 Reviewer Task Drawer；
  \item 支持批量撤回、批量转派。
\end{itemize}

\subsection{文案}

必须统一使用：

\begin{quote}
失败率 = 人工不合格数 / 抽检完成数
\end{quote}

不要再写成：

\begin{quote}
失败率 = 人工不合格数 / 批次总数
\end{quote}

\section{最终目标}

完成后，平台应形成完整闭环：

\begin{center}
采集批次 $\rightarrow$ 抽检质检 $\rightarrow$ 批次判定 $\rightarrow$ Episode 最终状态 $\rightarrow$ 任务级可用数据统计 $\rightarrow$ 合格数据清单导出
\end{center}

这将使平台从单纯质检工具升级为面向下游模型训练的数据资产供给系统。

\end{document}
